% ============================================================
% Capítulo: Resumen estructurado (Bloque B.2)
% ============================================================
\chapter*{Resumen}
\addcontentsline{toc}{chapter}{Resumen}

\textbf{PENDIENTE: se escribe al final, cuando el resto del documento
esté completo.} El resumen estructurado que exige el Bloque~B.2 de la
guía (problema, método, resultados principales con cifras reales,
conclusión) tiene que sintetizar el documento completo -- Introducción,
Trabajos Relacionados, Diseño y Arquitectura, Implementación, Resultados,
Ingeniería de Requisitos, Amenazas a la Validez, Discusión y
Conclusiones -- y varios de esos capítulos todavía no existen (ver la
nota de estado en \texttt{docs/informe-final.tex}). Redactarlo ahora
obligaría a inventar resultados y conclusiones que este documento
todavía no sostiene, exactamente el tipo de contenido fabricado que este
proyecto evita en el resto de su documentación.

\iffalse
Esqueleto de referencia para cuando se redacte (no es contenido real,
solo la estructura esperada por el Bloque B.2 -- eliminar todo este
bloque oculto al escribir el resumen definitivo):

- Problema/contexto: gestión bibliotecaria universitaria manual o
  parcialmente digitalizada, sin trazabilidad de requisitos ni evidencia
  empírica reproducible de calidad.
- Método: desarrollo del sistema SGB-SaaS (Spring Boot + Angular +
  PostgreSQL + Redis), con evaluación empírica en 4 frentes (rendimiento
  vía k6, seguridad vía auditoría OWASP Top 10, usabilidad vía SUS,
  ingeniería de requisitos vía checklist INCOSE v4), toda la evidencia
  versionada y reproducible.
- Resultados principales (con cifras reales, citando el archivo fuente
  de cada una -- no resumir de memoria): p95 de latencia con/sin cache,
  hallazgos OWASP corregidos, puntaje SUS (cuando exista), % de
  requisitos Must verificados, tasa de estabilidad de requisitos.
- Conclusión: qué tan bien se cumplieron los objetivos específicos
  (trazado a las RQ de la Introducción), y qué queda como limitación
  reconocida (no oculta).
\fi
